\chapter{Εισαγωγή}
\label{ch:intro}

\section{Το πρόβλημα}
\label{sec:intro-problem}

Η παρακολούθηση της δομικής ακεραιότητας (\en{Structural Health Monitoring}, SHM) γεφυρών, φραγμάτων, σηράγγων, πρανών και κτιρίων βασίζεται σε δίκτυα αισθητήρων που καταγράφουν συνεχώς μετακινήσεις, κλίσεις, ανοίγματα ρωγμών, πιέσεις και άλλα μεγέθη \cite{farrar2007shm}. Σε ένα τυπικό έργο, δεκάδες ή εκατοντάδες σημεία μέτρησης (κλισιόμετρα, ρωγμόμετρα, πιεζόμετρα, τοπογραφικοί στόχοι) τροφοδοτούν μια πλατφόρμα παρακολούθησης στο διαδίκτυο. Εκεί οι μετρήσεις αποθηκεύονται ως χρονοσειρές, συγκρίνονται με όρια συναγερμού και παρουσιάζονται ως πίνακες, διαγράμματα και κατόψεις.

Η πληροφορία όμως αφορά συγκεκριμένα σημεία του φυσικού χώρου, ενώ η παρουσίασή της γίνεται μακριά από αυτά, σε μια οθόνη. Ο μηχανικός που επισκέπτεται το έργο πρέπει να κάνει νοερά την αντιστοίχιση: ο αισθητήρας «Κ7» της λίστας, που η πλατφόρμα δείχνει σε κατάσταση συναγερμού, βρίσκεται σε κάποιο σημείο της κατόψεως, και το σημείο αυτό σε κάποια θέση του πραγματικού τοίχου, βάθρου ή πρανούς μπροστά του. Η αντιστοίχιση αυτή κοστίζει χρόνο, γίνεται συχνά λάθος σε έργα με πολλούς παρόμοιους αισθητήρες, και αφήνει αναξιοποίητη τη χωρική πληροφορία που έχει ο ίδιος ο χώρος: ποιοι αισθητήρες κινούνται μαζί, αν μια ζώνη μετακίνησης συμπίπτει με μια ορατή ρωγμή, πού βρίσκεται ο επόμενος αισθητήρας προς επιθεώρηση.

\review{Η παράγραφος περιγράφει τη ροή εργασίας «από γενική εμπειρία». Πρόσθεσε ένα \textbf{συγκεκριμένο παράδειγμα από δικό σου έργο ή επίσκεψη} (π.χ. το έργο όπου δοκιμάστηκε η εφαρμογή, πόσοι αισθητήρες, πόσος χρόνος χρειάζεται σήμερα για μια επιθεώρηση). Αν δεν υπάρχει μέτρηση, μην βάλεις αριθμούς.}

\section{Κίνητρο}
\label{sec:intro-motivation}

Η επαυξημένη και η μικτή πραγματικότητα (AR/MR) προσφέρουν μια άμεση απάντηση στο πρόβλημα: να εμφανίζεται η πληροφορία \emph{πάνω} στο σημείο που αφορά, στο οπτικό πεδίο του χρήστη \cite{azuma1997survey,bressa2022situated}. Η ιδέα δεν είναι νέα, αλλά μέχρι πρόσφατα προσέκρουε σε πρακτικά εμπόδια: ακριβές συσκευές οπτικής διαφάνειας με μικρό οπτικό πεδίο, ανάγκη για υπολογιστή ή καλώδια, και κυρίως η δυσκολία να τοποθετηθεί με ακρίβεια εκατοστών ένα γεωαναφερμένο σημείο στον χώρο μιας συσκευής που δεν γνωρίζει το γεωδαιτικό σύστημα αναφοράς.

Τα αυτόνομα headsets της τελευταίας γενιάς, όπως το Meta Quest~3, αλλάζουν την εικόνα. Συνδυάζουν έγχρωμο passthrough υψηλής ανάλυσης, εντοπισμό θέσης χωρίς εξωτερικό εξοπλισμό, αισθητήρα βάθους και spatial anchors που επιβιώνουν μεταξύ συνεδριών, με κόστος συγκρίσιμο με ένα tablet. Αυτό κάνει ρεαλιστική μια εφαρμογή πεδίου που ο μηχανικός μπορεί να φορέσει στο εργοτάξιο. Ανοιχτά παραμένουν τρία ερωτήματα: πόσο ακριβής μπορεί να είναι η ευθυγράμμιση με απλά μέσα, πόσο σταθερή μένει, και αν η εμπειρία βοηθά πραγματικά τον χρήστη.

\section{Ερευνητικά ερωτήματα}
\label{sec:rq}

Η εργασία εξετάζει τα εξής ερωτήματα:

\begin{description}
  \item[ΕΕ1 (ακρίβεια).] Με ποια ακρίβεια μπορεί ένα αυτόνομο headset μικτής πραγματικότητας να τοποθετήσει γεωαναφερμένα σημεία στη θέση τους στον πραγματικό χώρο, με λίγα σημεία ελέγχου και χωρίς εξωτερικό εξοπλισμό εντοπισμού;
  \item[ΕΕ2 (σταθερότητα).] Πόσο σταθερή μένει η ευθυγράμμιση στον χρόνο και μεταξύ συνεδριών μέσω spatial anchors, και πώς επηρεάζεται από τον φωτισμό και τις αλλαγές στο εργοτάξιο;
  \item[ΕΕ3 (απόδοση).] Ποια είναι τα όρια φόρτου (τρίγωνα μοντέλου, πλήθος αισθητήρων) για σταθερό ρυθμό καρέ σε αυτόνομη συσκευή, και πόσο τα μετακινούν οι τεχνικές βελτιστοποίησης;
  \item[ΕΕ4 (χρηστικότητα).] Βελτιώνει η επιτόπια απεικόνιση την ταχύτητα και την ακρίβεια μιας επιθεώρησης σε σχέση με τη συμβατική ροή εργασίας;
\end{description}

\section{Συνεισφορές}
\label{sec:contributions}

\begin{enumerate}
  \item Σχεδίαση και υλοποίηση ενός end-to-end συστήματος: πλατφόρμα παρακολούθησης (Terraweb) → εργαλεία PC σε Python → ενσωματωμένος εξυπηρετητής HTTP στη συσκευή → απεικόνιση σε μικτή πραγματικότητα (\cref{ch:design,ch:implementation}).
  \item Μέθοδος ευθυγράμμισης γεωαναφερμένου έργου με τον χώρο της συσκευής, με κλειστή λύση ελαχίστων τετραγώνων και πρόβλεψη σφάλματος σε πραγματικό χρόνο που καθοδηγεί τον χρήστη στο επόμενο σημείο ελέγχου (\cref{sec:impl-alignment}).
  \item Τεχνικές απεικόνισης δεδομένων αισθητήρων στον χώρο: κορδέλες ιστορικού, χρονομηχανή, διανύσματα μετακίνησης, χάρτης θερμότητας με παρεμβολή πάνω σε τριγωνοποίηση Delaunay, ράβδοι εγγύτητας στα όρια, ζωντανοί συναγερμοί, διαδρομή επιθεώρησης (\cref{sec:impl-vis}).
  \item Αλυσίδα προετοιμασίας και φόρτωσης βαριών μοντέλων (φωτογραμμετρία, DXF, IFC) σε συσκευή περιορισμένων πόρων (\cref{sec:impl-models}).
  \item Πειραματική αξιολόγηση ακρίβειας, σταθερότητας, απόδοσης και χρηστικότητας στο πεδίο (\cref{ch:evaluation}).
\end{enumerate}

\review[!]{Τα ΕΕ και οι συνεισφορές πρέπει στο τέλος να \textbf{ταιριάζουν με όσα δείχνει πράγματι το κεφ. 6}. Αν ένα πείραμα δεν γίνει (π.χ. η μελέτη χρηστών), αφαίρεσε ή μαλάκωσε το αντίστοιχο ΕΕ και τη συνεισφορά 5. Συζήτησέ τα με τον επιβλέποντα πριν ξεκινήσουν οι μετρήσεις.}

\section{Δομή της εργασίας}
\label{sec:outline}

Το \cref{ch:background} παρουσιάζει το θεωρητικό υπόβαθρο: μικτή πραγματικότητα, συστήματα συντεταγμένων, χωρική ευθυγράμμιση, παρεμβολή και απόδοση γραφικών σε κινητές συσκευές. Το \cref{ch:related} εξετάζει σχετικές εργασίες. Το \cref{ch:design} περιγράφει τις απαιτήσεις και την αρχιτεκτονική του συστήματος, και το \cref{ch:implementation} την υλοποίηση των υποσυστημάτων του. Το \cref{ch:evaluation} παρουσιάζει την πειραματική αξιολόγηση, το \cref{ch:discussion} συζητά τα αποτελέσματα και τους περιορισμούς, και το \cref{ch:conclusions} κλείνει με συμπεράσματα και μελλοντική εργασία.
